\documentclass[12pt]{article}
\usepackage{amsmath}
\usepackage{booktabs}
\usepackage{geometry}
\usepackage{hyperref}
\usepackage{graphicx}
\geometry{margin=1in}

\title{Problem Set 7\\Econ 5253 --- Spring 2026}
\author{Sumedha Ray}
\date{March 31, 2026}

\begin{document}
\maketitle

\section*{Question 6: Summary Statistics and Missing Data}

\input{summary_table}

The \texttt{logwage} variable is missing for 25.12\% of observations.
I believe \texttt{logwage} is most likely \textbf{Missing Not at Random (MNAR)}.
Women with very low or very high wages may be systematically less likely to
report earnings, meaning the probability of missingness depends on the unobserved
wage value itself. This violates both MCAR and MAR.

\section*{Question 7: Imputation Methods and Returns to Schooling}

I estimate the following model under four imputation strategies:
\begin{equation}
\log(\text{wage}_i) = \beta_0 + \beta_1 \text{hgc}_i + \beta_2 \text{college}_i
+ \beta_3 \text{tenure}_i + \beta_4 \text{tenure}_i^2
+ \beta_5 \text{age}_i + \beta_6 \text{married}_i + \varepsilon_i
\end{equation}

The true value of $\hat{\beta}_1 = 0.093$.

\input{regression_table}

All four methods underestimate the true $\hat{\beta}_1 = 0.093$, but the
degree of attenuation differs across methods. Mean imputation performs worst
($\hat{\beta}_1 = 0.0497$), which is expected: replacing missing wages with
the sample mean removes all covariance between log wages and the regressors,
attenuating every slope coefficient toward zero. Complete cases does better
($\hat{\beta}_1 = 0.0624$), though it relies on the MCAR assumption which is
likely violated here. Predicted value imputation ($\hat{\beta}_1 = 0.0571$)
is derived directly from the complete cases model, so it mechanically
preserves the covariate relationships but inflates precision by treating
imputed values as observed.

The MICE estimate ($\hat{\beta}_1 = 0.0583$) is the most principled of the
four. By drawing multiple plausible values for each missing observation via
predictive mean matching and pooling across five imputed datasets, MICE
correctly propagates uncertainty from the missing data into the standard
errors. That said, all estimates remain below the true value of 0.093,
suggesting the data are likely MNAR and that no imputation method fully
recovers the true parameter without modeling the missingness mechanism
directly.

\section*{Question 8: Project Progress}

My project, which I am developing as a potential job market paper,
asks whether economists who study sensitive topics---such as race, religion,
and discrimination---face systematically worse labor market and career
outcomes than observationally comparable candidates working on conventional
research topics. The core identification leverages variation in dissertation
topic (classified using JEL codes and keyword matching from ProQuest
Dissertations \& Theses) to ask three related questions: (1) does working
on sensitive topics reduce the probability of obtaining an academic
position, conditional on candidate quality and program rank? (2) do
researchers on these topics shift their agendas around the tenure
threshold, consistent with pre-tenure conformity pressure or post-tenure
intellectual liberation? and (3) to what extent do observed patterns
reflect demand-side discrimination versus rational supply-side
self-censorship? My advisor has suggested framing the topic selection
as a discrete choice/Roy model to account for selection into sensitive
research areas based on demographics and program characteristics. I am
currently building the data pipeline using ProQuest, RePec, and department
placement pages.
\end{document}
